% Môn: Vật lý
% Lớp: 11
% Chương: Chương III. Điện trường
% Bài: L11C3  Bài 16. Lực tương tác giữa hai điện tích
% Số câu: 9
% App đọc file này trực tiếp — sửa rồi Commit trên GitHub, bấm Đồng bộ GitHub .tex

% L11C3  Bài 16. Lực tương tác giữa hai điện tích

% ===== Câu 1 =====
% ID: L11C3B16-01-TL
% Mức: NB
\dangbt{Chưa có dạng}
\begin{ex}
% ID: L11C3B16-01-TL
% Mức: NB
Giải thích tại sao bụi bám chặt vào các cánh quạt máy bằng nhựa mặc dù các cánh quạt này thường xuyên quay rât nhanh. $\nguon{ly11 kntt.pdf}$
\choiceTF

\loigiai{
Do khi quay các cánh quạt cọ xát vào không khí nên bị nhiễm điện và hút các hạt bụi nhẹ trong không khí, làm chúng dính chặt vào cánh quạt.
}
\end{ex}

% ===== Câu 2 =====
% ID: L11C3B16-01-TN
% Mức: NB
\begin{ex}
% ID: L11C3B16-01-TN
% Mức: NB
Dùng vải cọ xát một đầu thanh nhựa rồi đưa lại gần hai vật nhẹ thì thấy thanh nhựa hút cả hai vật này. Hai vật này không thể là $\nguon{ly11 kntt.pdf}$
\choice
{hai vật không nhiễm điện.}
{hai vật nhiễm điện cùng loại.}
{\True hai vật nhiễm điện khác loại.}
{một vật nhiễm điện, một vật không nhiễm điện.}
\loigiai{
Chọn C.

Thanh nhựa hút được cả 2 vật chứng tỏ cả hai vật không thể nhiễm điện khác loại.
}
\end{ex}

% ===== Câu 3 =====
% ID: L11C3B16-02-TL
% Mức: TH
\begin{bt}
% ID: L11C3B16-02-TL
% Mức: TH
a) Tính lực tĩnh điện tương tác giữa hạt nhân nguyên tử helium với electron nằm trong lớp vỏ của nguyên tử này. Biết khoảng cách từ electron đến hạt nhân của nguyên tử helium là 2, 94 $\cdot$ 10^{-11} m, điện tích của electron là -1,6 $\cdot$  10^{-19} C. b) Nếu coi electron chuyển động tròn đều dưới tác dụng của lực hút tĩnh điện với bán kính quỹ đạo đã cho ở trên thì tốc độ góc và tốc độ của nó bằng bao nhiêu? Biết khối lượng của electron là 9,1 $\cdot$  10^{-31} $\mathrm{kg}$. $\nguon{ly11 kntt.pdf}$
\loigiai{
Chọn D.

Lực tương tác tỉ lệ nghịch với bình phương khoảng cách giữa hai điện tích.
}
\end{bt}

% ===== Câu 4 =====
% ID: L11C3B16-02-TN
% Mức: NB
\begin{ex}
% ID: L11C3B16-02-TN
% Mức: NB
Ba điện tích điểm chỉ có thể nằm cân bằng dưới tác dụng của các lực điện khi $\nguon{ly11 kntt.pdf}$
\choice
{ba điện tích cùng loại nằm ở ba đỉnh của một tam giác đều.}
{ba điện tích không cùng loại nằm ở ba đỉnh của một tam giác đều.}
{\True ba điện tích không cùng loại nằm trên cùng một đường thẳng.}
{ba điện tích cùng loại nằm trên cùng một đường thẳng.}
\loigiai{
Chọn C.

Ba điện tích điểm chỉ có thể nằm cân bằng dưới tác dụng của các lực điện khi ba điện tích không cùng loại nằm trên cùng một đường thẳng.
}
\end{ex}

% ===== Câu 5 =====
% ID: L11C3B16-03-TL
% Mức: TH
\begin{bt}
% ID: L11C3B16-03-TL
% Mức: TH
Hai quả cầu kim loại nhỏ có cùng kích thước, cùng khối lượng $90\,\mathrm{g}$, được treo vào cùng một điểm bằng hai sợi dây mảnh cách điện có cùng chiều dài $1,5\,\text{m}$. Truyền cho mỗi quả cầu một điện tích 2,4 \cdot  10^{-7} $C$ thì chúng đẩy nhau ra xa tới lúc cân bằng thì hai điện tích cách nhau một đoạn a. Coi góc lệch của hai sợi dây so với phương thẳng đứng là rất nhỏ. Tính độ lớn của a. Lấy $g =$ 10 $\mathrm{m/s}$ 2 $.$ \nguon{ly11kntt.pdf}
\loigiai{
Mỗi quả cầu chịu tác dụng của 3 lực: trọng lực $\overset\rightarrow{P}$; lực điện $\overset\rightarrow{F_{đ}}$ và lực căng $\overset\rightarrow{T}$.

Muốn quả cầu cân bằng phải có: $\overset\rightarrow{P}$ + $\overset\rightarrow{F_{đ}}$ + $\overset\rightarrow{T}$ = $\overset\rightarrow{0}$ hoặc $\overset\rightarrow{P}$ + $\overset\rightarrow{F_{đ}}$ = - $\overset\rightarrow{T}$, nghĩa là hợp lực của $\overset\rightarrow{P}$ và $\overset\rightarrow{F_{đ}}$ phải trực đối với $\overset\rightarrow{T}$ (Xem Hình $16.1\,\text{G}$).

[Hai quả cầu kim loại nhỏ có cùng kích thước]

Từ Hình 16.1 $G$ ta có: $\tan\alpha = \dfrac{F_{\text{d}}}{P} = \dfrac{F_{\text{đ}}}{mg} = \dfrac{9\cdot 10^{9}\dfrac{\left( {2,4\cdot 10^{-7}} \right)^{2}}{a^{2}}}{0,09.10}$ (1)

Vì góc α nhỏ nên ta có: $\tan\alpha = \sin\alpha = \dfrac{\text{a}}{2l} = \dfrac{a}{2.1,5}$ (2)

Từ (1) và (2) suy ra: a = $0,12\,\mathrm{m}$.
}
\end{bt}

% ===== Câu 6 =====
% ID: L11C3B16-03-TN
% Mức: NB
\begin{ex}
% ID: L11C3B16-03-TN
% Mức: NB
Tăng khoảng cách giữa hai điện tích lên 2 lần thì lực tương tác giữa chúng $\nguon{ly11 kntt.pdf}$
\choice
{tăng lên 2 lần.}
{giảm đi 2 lần.}
{tăng lên 4 lần.}
{\True giảm đi 4 lần.}
\loigiai{
Chọn D.

Lực tương tác tỉ lệ nghịch với bình phương khoảng cách giữa hai điện tích.
}
\end{ex}

% ===== Câu 7 =====
% ID: L11C3B16-04-TL
% Mức: TH
\begin{ex}
% ID: L11C3B16-04-TL
% Mức: TH
Một hệ gồm ba điện tích điểm dương q giống nhau và một điện tích điểm $Q$ nằm cân bằng. Biết ba điện tích q nằm ở ba đỉnh của một tam giác đều. Xác định dấu, độ lớn của điện tích (theo q) và vị trí của điện tích điểm Q. BÀl 17. KHÁ $I$ NIỆ $M$ ĐIỆ $N$ TRƯỜNG $\nguon{ly11 kntt.pdf}$
\choiceTF

\loigiai{
Thử xét trạng thái cân bằng của điện tích dương q đặt tại một trong ba đỉnh của tam giác đều ABC (cạnh a, đỉnh $C$ chẳng hạn). Lực đẩy của các điện tích q đặt tại hai đỉnh còn lại của tam giác lên điện tích đặt tại $C$ có độ lớn là: $F = \dfrac{q^{2}}{4\pi\varepsilon_{0}a^{2}}$

Hợp lực $\overset\rightarrow{F_{đ}}$ của hai lực này có phương nằm trên đường phân giác của góc $C$, chiều hướng ra ngoài tam giác (xem Hình $16.2\,\text{G}$), độ lớn: $F_{d} = F\sqrt{3} = \dfrac{q^{2}}{4\pi\varepsilon_{0}a^{2}}\sqrt{3}$ (1)

[Một hệ gồm ba điện tích điểm dương q giống nhau]

Muốn điện tích đặt tại $C$ nằm cân bằng thì phải có một lực hút cân bằng với lực đẩy $\overset\rightarrow{F_{đ}}$. Điện tích $Q$ do đó phải trái dấu với các điện tích q ($Q$ phải mang điện tích âm) và phải nằm trên đường phân giác của góc C. Tương tự, muốn cho các điện tích q đặt tại các đỉnh $A$ và $B$ nằm cân bằng thì điện tích $Q$ phải nằm trên các đường phân giác của góc $A$ và B. Nghĩa là $Q$ phải nằm tại trọng tâm của tam giác đều ABC và khoảng cách r từ $Q$ đến $C$ sẽ là: $r = \dfrac{2}{3}a \cdot \dfrac{\sqrt{3}}{2}$

Độ lớn của lực $\overset\rightarrow{F_{h}}$ do $Q$ tác dụng lên các điện tích q là: [Một hệ gồm ba điện tích điểm dương q giống nhau](2)

Vì F_(h) = F_(đ) nên từ (1) và (2), dễ dàng tính được độ lớn của $Q$ theo q: $Q$ = - $\dfrac{q}{\sqrt{3}}$.
}
\end{ex}

% ===== Câu 8 =====
% ID: L11C3B16-04-TN
% Mức: NB
\begin{ex}
% ID: L11C3B16-04-TN
% Mức: NB
Tăng đồng thời độ lớn của hai điện tích điểm và khoảng cách giữa chúng lên gấp đôi thì lực điện tác dụng giữa chúng $\nguon{ly11 kntt.pdf}$
\choice
{tăng lên 2 lần.}
{giảm đi 2 lần.}
{giảm đi 4 lần.}
{\True không đổi.}
\loigiai{
Chọn D.

Lực điện tỉ lệ thuận với độ lớn tích hai điện tích và tỉ lệ nghịch với bình phương khoảng cách giữa hai điện tích, khi tăng đồng thời độ lớn của hai điện tích điểm và khoảng cách giữa chúng lên gấp đôi thì lực điện tác dụng giữa chúng giữ nguyên.
}
\end{ex}

% ===== Câu 9 =====
% ID: L11C3B16-05-TN
% Mức: NB
\begin{ex}
% ID: L11C3B16-05-TN
% Mức: NB
Đối với điện trường xung quanh một điện tích điểm $Q$ đặt trong chân không, độ lớn của vectơ cường độ điện trường tại một điểm $M$ không phụ thộc vào $\nguon{ly11 kntt.pdf}$
\choice
{vị trí của điểm M.}
{\True dấu của điện tích Q.}
{độ lớn của điện tích Q.}
{khoảng cách từ điễm $M$ đến điện tích điểm Q.}
\loigiai{
Chọn B.

Độ lớn của cường độ điện trường không phụ thuộc vào dấu của điện tích.
}
\end{ex}
